\documentclass[11pt]{article}
\usepackage[letterpaper,landscape,margin=0.6in]{geometry}
\usepackage{amsmath}

\begin{document}
\pagestyle{empty}

\begin{center}
	{\LARGE\bfseries Ceramic Glaze Color-Response Test Program}
\end{center}
\hrule
\vspace{1.5em}

\noindent
\begin{minipage}[t]{0.46\textwidth}
	\textbf{Colorant Layout (one 5$\times$5 tile)} --- \% of dry base weight.

	\vspace{0.6em}
	\begin{tabular}{l c c c c c}
		\hline
		\textbf{Colorant}                    & \textbf{1} & \textbf{2} & \textbf{3} & \textbf{4} & \textbf{5} \\
		\hline
		Copper carbonate                     & 0.5        & 1          & 2          & 4          & 6          \\
		Red iron oxide ($\mathrm{Fe_2O_3}$)  & 1          & 2          & 4          & 8          & 12         \\
		Manganese dioxide ($\mathrm{MnO_2}$) & 1          & 2          & 4          & 8          & 12         \\
		Chrome oxide ($\mathrm{Cr_2O_3}$)    & 0.2        & 0.4        & 0.6        & 1          & 2          \\
		Cobalt oxide ($\mathrm{Co_3O_4}$)    & 0.2        & 0.4        & 0.6        & 1          & 2          \\
		\hline
	\end{tabular}
\end{minipage}
\hfill
\begin{minipage}[t]{0.50\textwidth}
	\textbf{Base Treatments (tiles 1--5)}

	\vspace{0.6em}
	\begin{tabular}{c l}
		\hline
		\textbf{Tile} & \textbf{Base treatment (every spot)} \\
		\hline
		1             & Base glaze (unmodified)              \\
		2             & Base + 5\% zircopax                  \\
		3             & Base + 10\% zircopax                 \\
		4             & Base + 4\% rutile                    \\
		5             & Base + 5\% zircopax + 4\% rutile     \\
		\hline
	\end{tabular}
\end{minipage}

\vspace{2em}

%% ============ Triaxial tests ============
\noindent\textbf{Triaxial Tests (tiles 6--10)} --- the 10 unique 3-colorant blends
($\binom{5}{3}=10$), two per tile. Corners: A = top-left, B = top-right,
C = bottom-left, D = bottom-right. The shared line blend runs A$\rightarrow$D, so the
two triaxials on each tile are A-B-D and A-C-D (corners = level-5 loading).

\vspace{0.6em}
\noindent\begin{tabular}{c l l l l}
	\hline
	\textbf{Tile} & \textbf{A}     & \textbf{B}     & \textbf{C} & \textbf{D}     \\
	\hline
	6             & Copper         & Manganese      & Chrome     & Red iron oxide \\
	7             & Copper         & Red iron oxide & Manganese  & Cobalt         \\
	8             & Copper         & Manganese      & Cobalt     & Chrome         \\
	9             & Red iron oxide & Chrome         & Cobalt     & Manganese      \\
	10            & Chrome         & Red iron oxide & Manganese  & Cobalt         \\
	\hline
\end{tabular}

\end{document}
